% id: 0981
% book: RPM 중학 3-1 · 08 이차함수의 그래프 ⑴
% section: 유형UP 21 이차함수의 그래프의 활용
% figure: crop:fig-0981.png
% answer: ④
% answer_source: 답지
% variant_level: 0 (원본 전사)
% 이 파일은 build.mjs 가 items.json 에서 만든다. 고칠 때는 items.json 을 고친다.
\smash{\raisebox{1.5mm}{\dmkichul{RPM 중학 3-1 0981번 · 140쪽}}}
\noindent\begin{minipage}[t]{0.58\linewidth}\vspace{0pt}\raggedright 오른쪽 그림과 같이 직선 $y=4$가 이차함수 $y=ax^2$의 그래프와 만나는 두 점을 $\pt{A}$, $\pt{E}$, 이차함수 $y=x^2$의 그래프와 만나는 두 점을 $\pt{B}$, $\pt{D}$, $y$축과 만나는 점을 $\pt{C}$라 하자. $\seg{AB}=\seg{BC}=\seg{CD}=\seg{DE}$일 때, 상수 $a$의 값은? \cond{두 점~$\pt{D}$, $\pt{E}$는 제$1$사분면 위의 점이다.}\end{minipage}\hfill\begin{minipage}[t]{0.4\linewidth}\vspace{0pt}\centering \includegraphics[width=98.6pt]{figures/fig-0981.png}\end{minipage}\par
\begin{choices32}{\rule[-3ex]{0pt}{8ex}$\dfrac{1}{16}$}{\rule[-3ex]{0pt}{8ex}$\dfrac{1}{12}$}{\rule[-3ex]{0pt}{8ex}$\dfrac{1}{8}$}{\rule[-3ex]{0pt}{8ex}$\dfrac{1}{4}$}{\rule[-3ex]{0pt}{8ex}$\dfrac{1}{2}$}\end{choices32}
